\section{Systematic Review Methodology}\label{sec:methodology}

We follow the PRISMA (Preferred Reporting Items for Systematic Reviews and Meta-Analyses) guidelines adapted for computational research \citep{Page2021}. This section details our search strategy, selection criteria, and data extraction protocol.

\subsection{Search Strategy}

We conducted systematic searches across five databases:
\begin{itemize}[leftmargin=*]
    \item \textbf{Scopus} --- comprehensive coverage of finance and CS journals
    \item \textbf{Web of Science} --- primary citation database for impact tracking
    \item \textbf{SSRN} --- working papers and preprints in finance
    \item \textbf{arXiv (q-fin, cs.LG, stat.ML)} --- preprints in quantitative finance and ML
    \item \textbf{Google Scholar} --- supplementary search for grey literature and conference proceedings
\end{itemize}

Search queries combined financial terms (\texttt{``stock prediction'' OR ``portfolio optimization'' OR ``risk management'' OR ``options pricing'' OR ``market microstructure'' OR ``quantitative finance''}) with ML terms (\texttt{``machine learning'' OR ``deep learning'' OR ``neural network'' OR ``random forest'' OR ``gradient boosting'' OR ``reinforcement learning'' OR ``natural language processing'' OR ``transformer'' OR ``LSTM''}).

Initial searches were conducted in March 2026, with supplementary forward-citation searches completed in April 2026.

\subsection{Inclusion and Exclusion Criteria}

\begin{table}[H]
\centering
\caption{Inclusion and exclusion criteria}
\label{tab:criteria}
\begin{tabularx}{\textwidth}{lX}
\toprule
\textbf{Include} & \textbf{Exclude} \\
\midrule
Published 2015--2025 & Published before 2015 \\
Applies ML to a quantitative finance problem & Pure ML methodology with no financial application \\
Presents empirical results on financial data & Theoretical/conceptual only \\
Journal article, conference paper, or working paper & Theses, textbooks, blog posts \\
English language & Non-English \\
Uses real or realistic simulated financial data & Uses only toy/synthetic examples \\
\bottomrule
\end{tabularx}
\end{table}

\subsection{Selection Process}

\begin{enumerate}[leftmargin=*]
    \item \textbf{Identification:} Database searches returned 4,783 unique records after deduplication. The arXiv q-fin section alone contributed over 12,600 submissions during our review period (2015--2025), growing from 852 total submissions in 2015 to 2,435 in 2024---a nearly threefold increase reflecting the rapid growth of quantitative finance research \citep{arXivStats2025}. Scopus, Web of Science, and SSRN contributed the remainder.
    \item \textbf{Screening:} Title and abstract screening against inclusion criteria reduced the set to 847 papers.
    \item \textbf{Eligibility:} Full-text review of screened papers identified 196 eligible studies.
    \item \textbf{Inclusion:} Forward and backward citation searches added 31 additional papers, yielding a final sample of 227 studies.
\end{enumerate}

\subsection{Data Extraction Protocol}

For each included study, we extracted the following fields into a structured database:

\begin{table}[H]
\centering
\caption{Data extraction fields}
\label{tab:extraction}
\begin{tabularx}{\textwidth}{lX}
\toprule
\textbf{Field} & \textbf{Description} \\
\midrule
Authors \& Year & Standard citation information \\
Venue & Journal, conference, or preprint server \\
Application Domain & Return prediction, risk management, portfolio optimization, derivatives, microstructure, or alternative data \\
ML Method(s) & Primary and benchmark methods used \\
Asset Class & Equities, fixed income, FX, commodities, crypto, multi-asset \\
Data Source & Specific database or data provider \\
Sample Period & Start and end dates of empirical analysis \\
Geography & US, Europe, Asia, emerging markets, global \\
Primary Metric & Main performance measure reported \\
Key Finding & One-sentence summary of the main result \\
Code Available & Yes/No with URL if available \\
Data Available & Public / Proprietary / Mixed \\
\bottomrule
\end{tabularx}
\end{table}

\subsection{Quality Assessment}

We do not exclude studies based on quality scores, as methodological standards differ substantially between the ML and finance communities. Instead, we flag quality dimensions that readers should consider when interpreting results:

\begin{itemize}[leftmargin=*]
    \item \textbf{Out-of-sample testing:} Does the study use a proper temporal train/validation/test split, or does it employ cross-validation (which can introduce look-ahead bias in time series)?
    \item \textbf{Transaction costs:} Are trading strategies evaluated net of realistic transaction costs?
    \item \textbf{Benchmark comparison:} Is the ML model compared against appropriate traditional benchmarks (e.g., GARCH for volatility, Fama-French for returns)?
    \item \textbf{Statistical significance:} Are results tested for statistical significance, or are only point estimates reported?
    \item \textbf{Reproducibility:} Is code and/or data publicly available?
\end{itemize}

These dimensions are reported in our comparison tables (\Cref{sec:appendix_tables}) to help readers assess the strength of evidence for each finding.
